\documentclass[11pt]{article}
% Shared preamble for manuscripts-v2 papers. Load with % Shared preamble for manuscripts-v2 papers. Load with % Shared preamble for manuscripts-v2 papers. Load with \input{../common/preamble.tex}
% immediately after \documentclass[11pt]{article}.
\ifdefined\pdftrailerid\pdftrailerid{}\fi
\usepackage[T1]{fontenc}
\usepackage{lmodern}
\usepackage[letterpaper,margin=1.02in]{geometry}
\usepackage{amsmath,amssymb,amsthm,mathtools,booktabs,array,microtype,enumitem}
\usepackage{tikz}
\usepackage{xurl}
\usepackage[colorlinks=true,allcolors=bluegreen]{hyperref}
\definecolor{bluegreen}{RGB}{22,83,91}
\usepackage[capitalize,nameinlink]{cleveref}

\theoremstyle{plain}
\newtheorem{theorem}{Theorem}[section]
\newtheorem{lemma}[theorem]{Lemma}
\newtheorem{proposition}[theorem]{Proposition}
\newtheorem{corollary}[theorem]{Corollary}
\newtheorem{conjecture}[theorem]{Conjecture}
\newtheorem*{theorem*}{Theorem}
\newtheorem{maintheorem}{Theorem}
\renewcommand{\themaintheorem}{\Alph{maintheorem}}
\theoremstyle{definition}
\newtheorem{definition}[theorem]{Definition}
\newtheorem{example}[theorem]{Example}
\newtheorem{question}[theorem]{Question}
\theoremstyle{remark}
\newtheorem{remark}[theorem]{Remark}
\crefname{maintheorem}{Theorem}{Theorems}
\crefname{conjecture}{Conjecture}{Conjectures}
\crefname{question}{Question}{Questions}
\crefname{example}{Example}{Examples}

\newcommand{\R}{\mathbb R}
\newcommand{\Q}{\mathbb Q}
\newcommand{\Z}{\mathbb Z}
\newcommand{\N}{\mathbb N}
\newcommand{\T}{\mathbb T}
\newcommand{\C}{\mathbb C}
\DeclareMathOperator{\den}{den}
\newcommand{\e}{\mathrm e}

% Immutable Lean source links at research-checkpoint-18.
\newcommand{\leancommit}{e4809cf561573aef588e6ffc491a35cdb79fd9ec}
\newcommand{\leanbase}{https://github.com/eldernos/MathematicsResearch/blob/\leancommit}
% \lean{File}{line}{label}: link to NormalityChecker/File.lean#Lline at the pinned commit.
\newcommand{\lean}[3]{\href{\leanbase/NormalityChecker/#1.lean\#L#2}{#3}}
% \repo{path}{label}: link to any repository path at the pinned commit.
\newcommand{\repo}[2]{\href{https://github.com/eldernos/MathematicsResearch/tree/\leancommit/#1}{#2}}

% Evidence labels (see papers/WAVE2-CONTRACT.md).
\newcommand{\evidence}[1]{\textup{[\emph{#1}]}}

\setlength{\emergencystretch}{1.5em}
\raggedbottom

% immediately after \documentclass[11pt]{article}.
\ifdefined\pdftrailerid\pdftrailerid{}\fi
\usepackage[T1]{fontenc}
\usepackage{lmodern}
\usepackage[letterpaper,margin=1.02in]{geometry}
\usepackage{amsmath,amssymb,amsthm,mathtools,booktabs,array,microtype,enumitem}
\usepackage{tikz}
\usepackage{xurl}
\usepackage[colorlinks=true,allcolors=bluegreen]{hyperref}
\definecolor{bluegreen}{RGB}{22,83,91}
\usepackage[capitalize,nameinlink]{cleveref}

\theoremstyle{plain}
\newtheorem{theorem}{Theorem}[section]
\newtheorem{lemma}[theorem]{Lemma}
\newtheorem{proposition}[theorem]{Proposition}
\newtheorem{corollary}[theorem]{Corollary}
\newtheorem{conjecture}[theorem]{Conjecture}
\newtheorem*{theorem*}{Theorem}
\newtheorem{maintheorem}{Theorem}
\renewcommand{\themaintheorem}{\Alph{maintheorem}}
\theoremstyle{definition}
\newtheorem{definition}[theorem]{Definition}
\newtheorem{example}[theorem]{Example}
\newtheorem{question}[theorem]{Question}
\theoremstyle{remark}
\newtheorem{remark}[theorem]{Remark}
\crefname{maintheorem}{Theorem}{Theorems}
\crefname{conjecture}{Conjecture}{Conjectures}
\crefname{question}{Question}{Questions}
\crefname{example}{Example}{Examples}

\newcommand{\R}{\mathbb R}
\newcommand{\Q}{\mathbb Q}
\newcommand{\Z}{\mathbb Z}
\newcommand{\N}{\mathbb N}
\newcommand{\T}{\mathbb T}
\newcommand{\C}{\mathbb C}
\DeclareMathOperator{\den}{den}
\newcommand{\e}{\mathrm e}

% Immutable Lean source links at research-checkpoint-18.
\newcommand{\leancommit}{e4809cf561573aef588e6ffc491a35cdb79fd9ec}
\newcommand{\leanbase}{https://github.com/eldernos/MathematicsResearch/blob/\leancommit}
% \lean{File}{line}{label}: link to NormalityChecker/File.lean#Lline at the pinned commit.
\newcommand{\lean}[3]{\href{\leanbase/NormalityChecker/#1.lean\#L#2}{#3}}
% \repo{path}{label}: link to any repository path at the pinned commit.
\newcommand{\repo}[2]{\href{https://github.com/eldernos/MathematicsResearch/tree/\leancommit/#1}{#2}}

% Evidence labels (see papers/WAVE2-CONTRACT.md).
\newcommand{\evidence}[1]{\textup{[\emph{#1}]}}

\setlength{\emergencystretch}{1.5em}
\raggedbottom

% immediately after \documentclass[11pt]{article}.
\ifdefined\pdftrailerid\pdftrailerid{}\fi
\usepackage[T1]{fontenc}
\usepackage{lmodern}
\usepackage[letterpaper,margin=1.02in]{geometry}
\usepackage{amsmath,amssymb,amsthm,mathtools,booktabs,array,microtype,enumitem}
\usepackage{tikz}
\usepackage{xurl}
\usepackage[colorlinks=true,allcolors=bluegreen]{hyperref}
\definecolor{bluegreen}{RGB}{22,83,91}
\usepackage[capitalize,nameinlink]{cleveref}

\theoremstyle{plain}
\newtheorem{theorem}{Theorem}[section]
\newtheorem{lemma}[theorem]{Lemma}
\newtheorem{proposition}[theorem]{Proposition}
\newtheorem{corollary}[theorem]{Corollary}
\newtheorem{conjecture}[theorem]{Conjecture}
\newtheorem*{theorem*}{Theorem}
\newtheorem{maintheorem}{Theorem}
\renewcommand{\themaintheorem}{\Alph{maintheorem}}
\theoremstyle{definition}
\newtheorem{definition}[theorem]{Definition}
\newtheorem{example}[theorem]{Example}
\newtheorem{question}[theorem]{Question}
\theoremstyle{remark}
\newtheorem{remark}[theorem]{Remark}
\crefname{maintheorem}{Theorem}{Theorems}
\crefname{conjecture}{Conjecture}{Conjectures}
\crefname{question}{Question}{Questions}
\crefname{example}{Example}{Examples}

\newcommand{\R}{\mathbb R}
\newcommand{\Q}{\mathbb Q}
\newcommand{\Z}{\mathbb Z}
\newcommand{\N}{\mathbb N}
\newcommand{\T}{\mathbb T}
\newcommand{\C}{\mathbb C}
\DeclareMathOperator{\den}{den}
\newcommand{\e}{\mathrm e}

% Immutable Lean source links at research-checkpoint-18.
\newcommand{\leancommit}{e4809cf561573aef588e6ffc491a35cdb79fd9ec}
\newcommand{\leanbase}{https://github.com/eldernos/MathematicsResearch/blob/\leancommit}
% \lean{File}{line}{label}: link to NormalityChecker/File.lean#Lline at the pinned commit.
\newcommand{\lean}[3]{\href{\leanbase/NormalityChecker/#1.lean\#L#2}{#3}}
% \repo{path}{label}: link to any repository path at the pinned commit.
\newcommand{\repo}[2]{\href{https://github.com/eldernos/MathematicsResearch/tree/\leancommit/#1}{#2}}

% Evidence labels (see papers/WAVE2-CONTRACT.md).
\newcommand{\evidence}[1]{\textup{[\emph{#1}]}}

\setlength{\emergencystretch}{1.5em}
\raggedbottom

% Author information approved by the repository owner on 2026-11-25:
%   name: Christopher Anderson; designation: Independent Researcher.
% No academic affiliation, ORCID or personal contact details have been supplied.
% Do not add any that the owner has not approved.
\newcommand{\OwnerAuthor}{Christopher Anderson}
\newcommand{\OwnerAffiliation}{Independent Researcher}
\newcommand{\OwnerContact}{Corrections and comments: \url{https://github.com/eldernos/MathematicsResearch/issues}}
\newcommand{\ReleaseLabel}{manuscripts-v2.0.0}
\newcommand{\LicenseLine}{\copyright\ 2026 Christopher Anderson. Licensed under \href{https://creativecommons.org/licenses/by/4.0/}{CC BY 4.0}.}
% \makebyline{version}{date}: centred author block used under the title.
\newcommand{\makebyline}[2]{%
  {\large \OwnerAuthor\par}%
  \ifx\OwnerAffiliation\empty\else{\small \OwnerAffiliation\par}\fi
  \vspace{0.3em}
  {\small Version #1 \enspace$\cdot$\enspace #2 \enspace$\cdot$\enspace \ReleaseLabel\par}%
  {\footnotesize \OwnerContact\par}%
  {\footnotesize \LicenseLine\par}%
  \hypersetup{pdfauthor={\OwnerAuthor}}}

\hypersetup{pdftitle={Exceptional Orbits and Discrepancy of Stoneham Products}}
\newcommand{\Disc}{\Delta}
\begin{document}
\begin{center}
{\LARGE\bfseries Exceptional Orbits and Discrepancy\\[2pt] of Stoneham Products\par}
\vspace{0.7em}
\makebyline{1.0}{25 November 2026}
\end{center}
\begin{abstract}
For the product of two Stoneham numbers with distinct odd prime parameters, we describe every subsequential distribution of the radix shifts as a mixture of Lebesgue measure and rational periodic orbits, and prove convergence of the scalar extreme interval discrepancy to that of the mixture. For the binary product $\alpha(2,3)\alpha(2,5)$, the upper interval discrepancy is at most $4/11+2.85\cdot10^{-10}$, and the exact bound $4/11$ is equivalent to the upper density of integer shifted prefixes being at most $4/11$. Ordinary binary normality remains open. A finite pair of local interpolation tests yields logarithmic equidistribution, lower discrepancy zero, and at most $O_\varepsilon(\log^* T)$ exceptional logarithmic scales up to $T$, with infinitely many admissible radices for every pair of distinct odd primes. We also record the kernel-verified $43/45$ bound for $\alpha(2,3)\alpha(2,5)$ together with transcendence and strict upper discrepancy for general multiplicatively independent integer parameters. Finally, for rational quadratics in a parameter-two Stoneham number, we combine the kernel-verified plateau law and resonant-radix classification with exact dyadic cycle discrepancies to obtain unconditional upper bounds strictly below one quarter.
\end{abstract}

\section{A discrepancy threshold}
For integers $B,c\ge2$, write
\[
 \alpha(B,c)=\sum_{j\ge1}\frac{B^{-c^j}}{c^j},\qquad
 \mu_N(x)=\frac1N\sum_{0\le n<N}\delta_{\{B^n x\}} .
\]
On the circle $\T=\R/\Z$ let $\lambda$ denote Lebesgue probability. For a Borel probability measure $\mu$ on $\T$, set
\[
 \Delta(\mu)=\sup_I|\mu(I)-|I||,
\]
where $I$ ranges over half-open arcs of $\T$. A wrapped arc has a nonwrapping complement with opposite centered mass, so the ordinary half-open interval convention in $[0,1)$ gives the same value, and $\Delta(\mu)=\sup_I(\mu(I)-|I|)$. Write $D_N(x)=\Delta(\mu_N(x))$, $D^+(x)=\limsup_{N\to\infty} D_N(x)$, and $D^-(x)=\liminf_{N\to\infty} D_N(x)$. Normality of $x$ in radix $B$ is equivalent to $D_N(x)\to0$: base-$B$ digit cylinders test all finite words and approximate every interval from inside and outside.

\begin{maintheorem}\label{th:binary}\evidence{Lean-supported paper proof}
Let $x=\alpha(2,3)\alpha(2,5)$. At the $n$th binary shift retain the terms already inserted and set
\[
 R_n=\sum_{\substack{i,j\ge1\\3^i+5^j\le n}}
       \frac{2^{n-3^i-5^j}}{3^i5^j},\qquad
 C_0=\limsup_{N\to\infty}\frac{\#\{0\le n<N:R_n\in\Z\}}{N}.
\]
Then $D^-(x)=0$, the orbit $\{2^n x\}$ is logarithmically equidistributed, and
\[
 D^+(x)\le \frac4{11}+2.85\cdot10^{-10},\qquad
 D^+(x)\le\frac4{11}\ \Longleftrightarrow\ C_0\le\frac4{11}.
\]
If either side of the equivalence fails, then $D^+(x)=C_0$.
\end{maintheorem}

The constant $4/11$ arises from the longest relative interval on which the rectangular array of summands lacks only its top corner. Such intervals occur arbitrarily late at every \emph{fixed} depth of simultaneous $3$-adic and $5$-adic cancellation. By contrast, an integer shifted prefix requires cancellation depths growing linearly with both shell indices. Ordinary binary normality of $\alpha(2,3)\alpha(2,5)$ remains open; the equivalence in Theorem~\ref{th:binary} isolates the zero-slack barrier $4/11$ as an integer-prefix density problem rather than proving normality. The proof of Theorem~\ref{th:binary} is completed in Section~\ref{sec:threshold}.

Stoneham introduced the numbers $\alpha(B,c)$ and proved normality for primitive-root radices by counting residues in rational expansions \cite{stoneham}; Bailey and Crandall established base-$B$ normality of $\alpha(B,c)$ for all coprime $B,c\ge2$ \cite[Corollary~4.9(i), p.~536]{bc}. Bailey and Borwein studied nonnormality in other bases and asked about sums of Stoneham constants \cite{bb}; see also Bugeaud's monograph \cite{bugeaud}. While rational affine transformations preserve normality \cite{wall} and additive normality preservers are characterized by zero entropy \cite{kamae,rauzy,weiss}, multiplication behaves quite differently: Bergelson and Downarowicz \cite{bd} construct counterexamples showing that deterministic factors need not preserve normality under multiplication. Recent public work has addressed sums and linear combinations of Stoneham series on distinct primes \cite{twosum,mira}. Multiplication of two Stoneham numbers on distinct primes introduces two-dimensional insertion locations $p^i+q^j$, whose reduced denominators depend on the boundary geometry of the cutoff. We first analyze the limiting measures, local $p$-adic interpolation, scale sparsity, and discrepancy threshold for independent-prime products, then record the kernel-verified $43/45$ bound and general independent-parameter transcendence theorem, and finally treat the analogous exceptional plateau cycles arising from rational quadratics on the dyadic axis $c=2$.

\section{The limiting measures}\label{sec:mixture}
Fix distinct odd primes $p<q$ coprime to $B$ and put $x=\alpha(B,p)\alpha(B,q)$. Define the shifted rational prefix $R_n$ as in Theorem~\ref{th:binary} with $(B,p,q)$ replacing $(2,3,5)$. For $n\ge q$, put $K=\lfloor\log_p n\rfloor$, $L=\lfloor\log_q n\rfloor$, $P=p^K$, and $Q=q^L$. Intersect each shell of constant $(K,L)$ with the half-open intervals between consecutive insertions (each containing its left insertion point). Call every nonempty interval $I=[U,V)$ obtained in this way a \emph{plateau}. Its unshifted prefix $B^{-U}R_U$ is constant on $I$, so $\{R_n\}$ follows the radix-$B$ cycle $c(I)$ of $\{R_U\}$, which is purely periodic because its reduced denominator $m_I=\den(R_U)$ divides $PQ\le n^2$ and is coprime to $B$. Apply the same partition on $[0,q)$ with empty prefix $R_n=0$.

Let $\nu_c$ denote the uniform probability measure on cycle $c$, indexing all rational $B$-cycles whose reduced denominator $m_c$ is supported on $\{p,q\}$ (including the singleton $\{0\}$ of denominator $1$). For each $N\ge1$ and each cycle $c$, put
\[
 w_{c,N}=\frac1N\sum_{c(I)=c}\bigl|I\cap[0,N)\bigr|,
\]
where $|S|$ denotes integer cardinality. Let $\mathcal W$ be the set of coordinatewise subsequential limits $w=(w_c)_c$ of $(w_{c,N})_c$ as $N\to\infty$. Every $w\in\mathcal W$ satisfies $w_c\ge0$ and $\sum_c w_c\le1$ by considering finite partial sums.

\begin{theorem}\label{th:mixture}\evidence{Lean-supported paper proof}
Along every sequence $N\to\infty$ on which $w_{c,N}\to w_c$ for each cycle $c$,
\[
 \mu_N(x)\ \longrightarrow\ \sigma(w)=
     \Bigl(1-\sum_c w_c\Bigr)\lambda+\sum_c w_c\nu_c
 \quad\text{weakly on }\T,\qquad
 D_N(x)\longrightarrow\Delta(\sigma(w)).
\]
In particular, $D^+(x)=\max_{w\in\mathcal W}\Delta(\sigma(w))$.
\end{theorem}
\begin{proof}
We use the \lean{SmoothOrbit}{207}{fixed-support orbit estimate}, a qualitative counterpart of Vandehey's short-orbit bound \cite[Corollary~1.2]{vandehey}. For fixed integer $d\ge1$, a reduced rational $y$ with denominator $m$ supported on $\{p,q\}$ (coprime to $B$) satisfies
\[
 \frac1T\sum_{i=0}^{T-1}\e(B^i y)=o_{m\to\infty}(1)\qquad (m\le T^d),
\]
uniformly in $T$ and the numerator of $y$, where $\e(t)=\exp(2\pi\mathrm i t)$. Multiplying $y$ by a fixed nonzero integer $h$ reduces the denominator by at most a factor $|h|$, so the same uniform bound holds for every fixed nonzero frequency $h$. Sandwiching the indicator of any arc with endpoints on a fixed finite mesh between continuous functions approximated uniformly by trigonometric polynomials, and enclosing arbitrary arcs between mesh arcs, shows that the uniform arc discrepancy of $\{B^i y\}_{0\le i<T}$ is $o_{m\to\infty}(1)$ uniformly for $m\le T^d$. Applying this to arbitrarily large multiples of one full period yields $\sup_{m_c>M}\Delta(\nu_c)\to0$ as $M\to\infty$.

Let $\widetilde\mu_N=\frac1N\sum_{0\le n<N}\delta_{\{R_n\}}$ be the empirical measure of the rational prefixes. Partition $[0,N)$ into blocks of length $L_N=\lfloor\sqrt N\rfloor$. Below $N$ there are at most $\log_p N+\log_q N+2$ shell boundaries and at most $(\log_p N)(\log_q N)$ insertions, so the final incomplete block and all blocks meeting a plateau boundary together occupy $O(N^{-1/2}(1+\log^2 N))$ of $[0,N)$. Every remaining clean block lies in a single plateau $I$ with $m_I\le N^2\le L_N^5$ for all large $N$. Given $\varepsilon>0$, choose $M$ so large that every clean block with $m_I>M$ has arc error at most $\varepsilon L_N$ uniformly over all arcs $J\subset\T$. On a clean block with $m_I\le M$, the orbit is periodic of period at most $m_I\le M$, so the count in any arc $J$ differs from $L_N\nu_{c(I)}(J)$ by at most $M$. Summing over at most $\lceil N/L_N\rceil$ blocks gives total incomplete-cycle error $O(M/\sqrt N)$. Defining the truncated mixture
\[
 \sigma_{N,M}=\Bigl(1-\sum_{m_c\le M}w_{c,N}\Bigr)\lambda+\sum_{m_c\le M}w_{c,N}\nu_c,
\]
the discarded blocks affect both $\widetilde\mu_N$ and the weights $w_{c,N}$ by at most $O(N^{-1/2}(1+\log^2 N))$, so
\[
 \sup_J\bigl|\widetilde\mu_N(J)-\sigma_{N,M}(J)\bigr|\le \varepsilon+O\!\left(\frac{M+\log^2 N}{\sqrt N}\right).
\]
Along a sequence where $w_{c,N}\to w_c$ for each $c$, the finite sum $\sigma_{N,M}$ converges uniformly on all arcs to $\sigma_M=(1-\sum_{m_c\le M}w_c)\lambda+\sum_{m_c\le M}w_c\nu_c$. Moreover,
\[
 \sup_J\bigl|\sigma(w)(J)-\sigma_M(J)\bigr|
 =\sup_J\Bigl|\sum_{m_c>M}w_c\bigl(\nu_c(J)-|J|\bigr)\Bigr|
 \le \sup_{m_c>M}\Delta(\nu_c).
\]
Letting $N\to\infty$ and then $M\to\infty$ (with $\varepsilon\to0$) proves uniform arc convergence $\sup_J|\widetilde\mu_N(J)-\sigma(w)(J)|\to0$, and hence $\Delta(\widetilde\mu_N)\to\Delta(\sigma(w))$.

To pass to the real shifts $\{B^n x\}$, observe that every omitted pair $(i,j)$ with $p^i+q^j>n$ satisfies $B^{n-p^i-q^j}\le 1$ and either $p^i>n/2$ or $q^j>n/2$. Summing the two geometric coefficient tails gives
\[
 0\le B^n x-R_n\le \frac{2(p+q)}{(p-1)(q-1)\,n}=O\!\left(\frac1n\right).
\]
For every $f\in C(\T)$, uniform continuity implies $\int f\,d\mu_N(x)-\int f\,d\widetilde\mu_N\to0$, while approximating $f$ uniformly by step functions on arcs yields $\int f\,d\widetilde\mu_N\to\int f\,d\sigma(w)$; thus $\mu_N(x)\to\sigma(w)$ weakly. Weak convergence alone does not control $\mu_N(x)(J)$ when an endpoint of $J$ is an atom of $\sigma(w)$, so we transfer the scalar discrepancy directly. Given $t>0$, choose $K$ so that $|B^n x-R_n|\le t$ for $n\ge K$. For any half-open arc $J$ and any $\eta>0$, let $J^+$ be the half-open arc extending $J$ by $t+\eta$ at both ends (or all of $\T$ if $|J|+2(t+\eta)\ge1$). Every $n\ge K$ with $\{B^n x\}\in J$ satisfies $\{R_n\}\in J^+$, giving
\[
 \mu_N(x)(J)-|J|\le \widetilde\mu_N(J^+)-|J^+|+2(t+\eta)+\frac KN
 \le \Delta(\widetilde\mu_N)+2(t+\eta)+\frac KN.
\]
Taking the supremum over $J$, letting $\eta\downarrow0$, and interchanging the two sequences yields $|D_N(x)-\Delta(\widetilde\mu_N)|\le 2t+K/N$; letting $N\to\infty$ and then $t\downarrow0$ gives $D_N(x)-\Delta(\widetilde\mu_N)\to0$ (the \lean{CircleDiscrepancy}{137}{circle perturbation theorem}). Combining this with $\Delta(\widetilde\mu_N)\to\Delta(\sigma(w))$ proves $D_N(x)\to\Delta(\sigma(w))$. Finally, since the set of cycles is countable, Cantor diagonalization on a sequence realizing $D^+(x)$ extracts a limit $w\in\mathcal W$ with $\Delta(\sigma(w))=D^+(x)$, proving that the maximum is attained.
\end{proof}

\section{Denominator losses and logarithmic scales}\label{sec:local}
For each \emph{ordered} pair $(p,q)$ of distinct odd primes coprime to $B$, use $p$ for the prime at which the denominator is tested. For $n\ge\max(p,q)$ (so $K,L\ge1$), define the shape coordinates
\[
 a=L-\#\{j\ge1: q^j\le n-P\},\qquad
 d=K-\#\{i\ge1: p^i\le n-Q\}.
\]
They are both zero on a full rectangle ($P+Q\le n$) and both positive otherwise. All cells with $i<K,j<L$ occur since $P/p+Q/q\le(1/3+1/5)n<n$. Exchanging $(p,q)$ exchanges $(a,d)$.

In $\Z_p$ put $x_j=q^{-j}B^{-q^j}$, $S_m=\sum_{j=1}^m x_j$, and $S_0=0$. For $a,d\ge1$ set
\[
 G_{a,d}(m)=(1-p)S_{m-a}+pS_{m-1}+p^d x_m ,
 \qquad G_{0,0}(m)=S_m .
\]
We say the ordered axis $(p,q)$ satisfies the \emph{simple-zero conditions} if there is an integer $u\ge1$ such that, with $\omega\in\Z_p^\times$ the Teichm\"uller representative of $B$,
\[
 q^u\equiv1\pmod p,\quad
 \omega^{q^{j+u}}=\omega^{q^j}\ (j\ge1),\quad
 p\nmid S_i\ (1\le i<u),\quad v_p(S_u)=1 .
\]
The torsion test is finite: check the eventual cycle and initial terms of $q^j\bmod\operatorname{ord}_p(B)$. At the reversed axis $(q,p)$ write $v$ for the corresponding period.

\begin{proposition}\label{th:local}\evidence{Lean-supported paper proof}
For every $n\ge\max(p,q)$, the number of $p$-layers lost from $p^K$ in $\den(R_n)$ equals $\min(K,v_p(G_{a,d}(L)))$. Under the simple-zero conditions, $G_{a,d}(m)$ is a $p$-adic unit off the sheet $m\equiv a\bmod u$. On that sheet there is a unique $\zeta_{p,a,d}\in\Z_p$ such that
\[
 v_p\bigl(G_{a,d}(a+uz)\bigr)=1+v_p(z-\zeta_{p,a,d}) \qquad (z\in\Z_p).
\]
For a full rectangle, $\zeta_{p,0,0}=0$. Moreover, for every pair of distinct odd primes $(p,q)$, infinitely many integer radices $B\ge2$ coprime to $pq$ satisfy the simple-zero conditions at both primes.
\end{proposition}
\begin{proof}
Put $A_n=PQR_n$ and $I_j=\#\{i\ge1:p^i\le n-q^j\}$. For $1\le i\le K$, $\varphi(p^i)=p^{i-1}(p-1)$ divides $P-p^i=p^i(p^{K-i}-1)$, so Euler's congruence gives $B^{P-p^i}\equiv1\bmod p^i$; multiplying by $p^{K-i}$ restores precision $p^K$. Summing $(p-1)\sum_{i=1}^{I_j}p^{K-i}=p^K-p^{K-I_j}$ yields
\[
 (p-1)B^P A_n\equiv -B^n Q\sum_{j=1}^L p^{K-I_j}x_j\pmod{p^K}.
\]
In a nonrectangle the fibre weights $p^{K-I_j}$ are $1$ for $1\le j\le L-a$, $p$ for $L-a<j<L$, and $p^d$ for $j=L$ (including the empty-fibre boundary cases $a=L$ where the first range is empty and $d=K$ where $p^K x_L\equiv0\bmod p^K$). Both outer factors $(p-1)B^P$ and $-B^n Q$ are $p$-adic units. This proves the loss identity; for the binary pair $(2,3,5)$ see the \lean{ProductThreeFiveInterpolation}{168}{ternary} and \lean{ProductThreeFiveInterpolation}{182}{quinary endpoint congruences} and the \lean{ProductThreeFiveInterpolation}{123}{shape identity}. Modulo $p$, $G_{a,d}(m)\equiv S_{m-a}$, which vanishes modulo $p$ precisely when $u\mid(m-a)$ by periodicity of $x_j\bmod p$ and the third condition.

Write $A=q^u$, $\ell=\log_p(B/\omega)\in p\Z_p$, and $g(z)=\sum_{i=1}^u x_{i+uz}$. On the sheet $i+uz$,
\[
 x_{i+uz}=q^{-i}\omega^{-q^i}A^{-z}\exp(-q^i\ell A^z).
\]
Since $\log_p A,\ell\in p\Z_p$ and $p$ is odd, these formulas extend analytically to $z\in\Z_p$. Let $\nabla$ denote the forward difference operator. Writing $\exp(-q^i\ell A^z)=\exp(-q^i\ell)\exp(-q^i\ell(A^z-1))$ separates $x_i$ and expands the remaining factor in powers of $A^z-1$:
\[
 x_{i+uz}=x_i\sum_{h=0}^\infty\frac{(-q^i\ell)^h}{h!}
          \sum_{j=0}^h(-1)^{h-j}\binom hj A^{(j-1)z}.
\]
Applying $\nabla^k|_{z=0}$ term by term gives $\nabla^k(A^{(j-1)z})|_{z=0}=(A^{j-1}-1)^k\in p^k\Z_p$. For $h=0$, summing over $1\le i\le u$ gives $S_u(A^{-1}-1)^k$, of valuation at least $k+1$. For $h\ge1$, $v_p((-q^i\ell)^h/h!)\ge h-(h-1)/(p-1)\ge1$ since $p\ge3$, tending to infinity with $h$. Hence $v_p(\nabla^k g(0))\ge k+1$ for all $k\ge0$. The Mahler--Newton series \cite{mahler,robert}
\[
 F(z)=\sum_{k\ge0}\nabla^k g(0)\binom z{k+1}
\]
converges uniformly on $\Z_p$, equals $S_{uz}$ at all $z\in\Z_{\ge0}$ by telescoping, and defines its continuous extension.

For $y\in\Z_p$ with $v_p(y)=t$, Vandermonde's identity and $\binom yj=(y/j)\binom{y-1}{j-1}$ give valuation at least $t+k+1-\lfloor\log_p(k+1)\rfloor\ge t+2$ for the increment of each term $k\ge1$ (since $p\ge3$ implies $k+1-\lfloor\log_p(k+1)\rfloor\ge2$ for all $k\ge1$). Thus $F(z+y)-F(z)\equiv yS_u\bmod p^{t+2}$. Every individual extra summand $x_{uz+j}$ changes by a multiple of $p^{t+1}$ because $A^y\equiv1\bmod p^{t+1}$. On $m=a+uz$, writing $G_{a,d}(a+uz)=S_{uz}+p\sum_{j=1}^{a-1}x_{uz+j}+p^d x_{uz+a}$ shows that $G_{a,d}$ obeys the same congruence modulo $p^{t+2}$. Therefore $z\mapsto G_{a,d}(a+uz)/p$ is an isometry of $\Z_p$, surjective because it induces a bijection on every finite quotient $\Z/p^r\Z$. Its unique zero $\zeta_{p,a,d}$ satisfies the stated valuation formula, valid also when $p\mid u$. For $(a,d)=(0,0)$, $G_{0,0}(0)=S_0=0$ gives $\zeta_{p,0,0}=0$.

Finally, impose $B\equiv1+pc\bmod p^2$, so $\omega=1$ and $x_j\equiv q^{-j}-pc\bmod p^2$. If $q\not\equiv1\bmod p$, take $u=\operatorname{ord}_p(q)>1$; then $p\nmid u$, the partial sums vanish modulo $p$ exactly at multiples of $u$, and only one class of $c\bmod p$ makes $S_u\equiv0\bmod p^2$, leaving $p-1\ge2$ valid classes. If $q\equiv1\bmod p$, take $u=p$; writing $q^{-1}=1+pr$ gives $S_i\equiv i\bmod p$ for $1\le i<p$ and $S_p\equiv p+pr\sum_{j=1}^p j\equiv p\bmod p^2$ for every $c$ because $p$ is odd. Combining valid classes at $p$ and $q$ by the Chinese remainder theorem modulo $(pq)^2$ yields infinitely many admissible integers $B\ge2$.
\end{proof}

\begin{proposition}\label{th:eisen}\evidence{Paper proof}
For $(B,p,q)=(2,3,5)$ or $(2,5,3)$ and every $a,d\ge1$, one has $\zeta_{p,a,d}\notin\Z$.
\end{proposition}
\begin{proof}
The periods are $u=2$ and $u=4$, respectively. The summands $x_j\bmod p$ repeat $(1,2)$ and $(4,2,1,3)$, with block sums $S_2\equiv6\bmod9$ and $S_4\equiv15\bmod25$, so Proposition~\ref{th:local} applies. Suppose $z\in\Z$ and put $m=a+uz$. If $m\ge a$, then $G_{a,d}(m)=S_{m-a}+p(S_{m-1}-S_{m-a})+p^d x_m$ is a strictly positive rational number and cannot vanish.

For $m<a$, extend $S_j$ to all $j\in\Z$ by $S_j-S_{j-1}=x_j$ with $S_0=0$, where $x_{i+uz}$ ($1\le i\le u$, $z\in\Z$) is given by the analytic formula of Proposition~\ref{th:local}. Since $F(z+1)-F(z)=g(z)$ holds on $\Z_{\ge0}$, continuity extends it to all $z\in\Z$, so this backward recurrence matches the continuous interpolation on every sheet. Note that $\omega^{q^{u(-z)}}=\omega$ for every $z\in\Z$: at $p=3$ we have $\omega=-1$ and $q=5$ is odd, while at $p=5$ we have $\omega^4=1$ and $3^4\equiv1\bmod4$. Raising $q^{i+uz}x_{i+uz}=\omega^{-q^i}\exp(-q^i\ell q^{uz})$ to the power $q^h$ when $i+uz=-h\le0$ therefore yields $\omega^{-q^{i+h}}\exp(-\ell)=\omega^{-1}(2/\omega)^{-1}=1/2$. In particular $x_0=1/2$, and for $h\ge1$ the factor $2^{-q^{-h}}$ in $x_{-h}=q^h 2^{-q^{-h}}$ is the unique $q^h$th root of $1/2$ in $\Z_p^\times$ (raising to $q$ is bijective on $\Z_p^\times$ since $\gcd(q,p(p-1))=1$ in both cases). These roots are compatible across $h$.

Using $S_{-j}=-\sum_{i=1-j}^0 x_i$ for $j\ge1$, expand $G_{a,d}(m)=(1-p)S_{m-a}+pS_{m-1}+p^d x_m$. If $m\le a-2$, a negative index appears; the smallest negative index is $-H=m-a+1\le-1$ with nonzero coefficient $p-1$ when $a>1$, and $-H=m\le-1$ with nonzero coefficient $p^d-1$ when $a=1$. Every negative-index term $x_i$ ($-H\le i\le-1$) is a rational multiple of $\theta^{q^{H+i}}$ where $\theta=2^{-1/q^H}$, with distinct exponents $1,q,\ldots,q^{H-1}<q^H$, while all terms with $i\ge0$ are rational. Because $X^{q^H}-2$ is Eisenstein at $2$, we have $[\Q(\theta):\Q]=q^H$, so $\{1,\theta,\ldots,\theta^{q^H-1}\}$ is $\Q$-linearly independent; since the coefficient of $\theta^1$ is nonzero, $G_{a,d}(m)\ne0$. In the remaining case $m=a-1\ge0$, if $m\ge1$ then $(1-p)S_{-1}=(p-1)/2$, $pS_{m-1}\ge0$, and $p^d x_m>0$, while at $m=0,a=1$ the value is $(p^d-1)/2>0$. Hence no integer $z$ is a zero.
\end{proof}

\begin{theorem}\label{th:sparse}\evidence{Lean-supported paper proof}
Suppose both ordered axes $(p,q)$ and $(q,p)$ satisfy the simple-zero conditions. Then:
\begin{enumerate}[label=\textup{(\roman*)},leftmargin=2em,itemsep=2pt]
\item For fixed integers $h,k\ge1$, the set of positions $n$ losing at least $h$ layers at $p$ and at least $k$ layers at $q$ has logarithmic density $1/(uvp^{h-1}q^{k-1})$, and the same conditional logarithmic density on every fixed shape $(a,d)$ of positive logarithmic density.
\item For every $\varepsilon>0$,
\[
 \#\bigl\{m\in\N:m\le T,\ \exists\,N\in[e^m,e^{m+1})\text{ with }D_N(x)>\varepsilon\bigr\}
       =O_{B,p,q,\varepsilon}(\log^*T),
\]
where $\log^*$ counts iterated natural logarithms until the argument is at most $1$. Consequently,
\[
 \sum_{\substack{N\ge2\\D_N(x)>\varepsilon}}\frac1{N\log N}<\infty,\qquad
 \sum_{\substack{2\le N\le X\\D_N(x)>\varepsilon}}\frac1N
       =O_{B,p,q,\varepsilon}(1+\log^*\log X).
\]
\item The radix orbit $\{B^n x\}$ is logarithmically equidistributed on $\T$, and $D^-(x)=0$.
\end{enumerate}
\end{theorem}
\begin{proof}
(i) For fixed depths $h,k\ge1$ and large shells, Proposition~\ref{th:local} reduces the simultaneous loss conditions on shape $(a,d)$ to one residue class of $L\bmod up^{h-1}$ and one residue class of $K\bmod vq^{k-1}$. Put $t=\log n$ and run the continuous linear flow with coordinates $t/\log p\bmod vq^{k-1}$ and $t/\log q\bmod up^{h-1}$. By unique factorization, $\log p/\log q$ is irrational, so every nontrivial character has nonzero frequency and time average zero; trigonometric approximation gives Haar equidistribution for sets with null boundary. The fractional parts $s=\{t/\log p\}$ and $t'=\{t/\log q\}$ determine the shape (for instance, $a=-\lfloor t'+\log_q(1-p^{-s})\rfloor$ on nonrectangles, clipped to $0$ on rectangles) up to empty fibres occurring within distance $q$ of a power $p^K$ or distance $p$ of a power $q^L$, whose total logarithmic weight is bounded. The shape discontinuities are axes and countably many level curves accumulating only at the axes, so truncating arbitrarily thin strips near the axes yields Riemann-integrable indicators. At each fractional pair, exactly one residue cell on each axis passes. Finally, physical memberships change only at powers or sums of powers, and replacing $\int_n^{n+1}dt/t=\log(1+1/n)$ by $1/n$ changes the sum by $O(1)$. This proves (i).

(ii) Fix a denominator cap $M\ge1$. On full rectangles, Proposition~\ref{th:local} gives $\zeta_{p,0,0}=0$, so the lost $p$-layers are at most $1+\lfloor\log_p L\rfloor$; thus $\den(R_n)\to\infty$ uniformly on full rectangles, leaving only finitely many rectangular positions with $\den(R_n)\le M$. For a fixed nonrectangular shape $(a,d)$, two distinct active shells $(K_1,L_1),(K_2,L_2)$ with losses at least $(h_1,k_1),(h_2,k_2)$ satisfy
\[
 up^{\min(h_1,h_2)-1}\mid(L_2-L_1),\qquad
 vq^{\min(k_1,k_2)-1}\mid(K_2-K_1),
\]
by subtracting the residue tests of Proposition~\ref{th:local} (for the binary case see the \lean{Phase16StonehamBounds}{291}{simultaneous ultrametric spacing inequality}). Overlapping shells satisfy $-\log q<L\log q-K\log p<\log p$, and the indices increase componentwise along plateaus. If one index stays fixed, the change of the other is bounded by a constant depending only on $p,q$, which is eventually smaller than the dividing modulus once $\den(R_n)\le M$. Hence for distinct sufficiently late shells of shape $(a,d)$ with $\den(R_n)\le M$, ordering $U_1<U_2$ forces $K_1<K_2$ and $L_1<L_2$. Since $h_i\ge K_i-\log_p M$ and $k_i\ge L_i-\log_q M$, we have $\min(h_1,h_2)\ge K_1-\log_p M$, so $L_2>L_2-L_1\ge up^{K_1-\lfloor\log_p M\rfloor-1}\gg_{B,p,q,M} U_1$, giving $\log U_2\gg_{B,p,q,M} U_1$. (More generally, retaining at most fractions $\theta<1$ of both exponents gives $\log U_2\gg U_1^{1-\theta}$.)

For context, on the torsion classes $L\equiv a\bmod2$, $K\equiv d\bmod4$ of the binary product $(2,3,5)$, bounded denominators $3^r5^s$ at $K>r,L>s$ are equivalent to simultaneous approximation
\[
 v_3(L-a-2\zeta_{3,a,d})\ge K-r-1,\qquad
 v_5(K-d-4\zeta_{5,d,a})\ge L-s-1 .
\]
If either root were algebraic over $\Q$, a nonzero integral annihilating polynomial $f$ would give $v_p(m-\alpha)\le v_p(f(m))=O_f(\log m)$ outside its finite integer zero set, ruling out such shells; Proposition~\ref{th:eisen} rules out integer roots but does not supply an algebraic height bound.

To control all shapes simultaneously, for $0<\eta<1/2$ discard the boundary strips where $P>(1-\eta)n$ or $Q>(1-\eta)n$. On the $p$-axis, setting $t=p^{\lfloor\log_p(N-1)\rfloor}/N\in(1/p,1]$, the last strip contributes $\min(1-t,t\eta/(1-\eta))$ and earlier strips contribute $t\eta/((p-1)(1-\eta))+O((\log N)/N)$; this piecewise-linear function peaks at $t=1-\eta$ with upper density $p\eta/(p-1)$ (see the \lean{Phase16StonehamBounds}{248}{boundary-strip bound}). Outside both strips, every nonrectangle with $K,L\ge1$ satisfies $q^{L-a+1}>n-P\ge\eta n\ge\eta Q$, hence $q^{a-1}<1/\eta$, and symmetrically $p^{d-1}<1/\eta$, leaving finitely many bulk shapes.

Given $\varepsilon>0$, the block proof of Theorem~\ref{th:mixture} yields $M$ such that $D_N(x)>\varepsilon$ eventually requires $\#\{n<N:\den(R_n)\le M\}>\varepsilon N/2$. Choose $\eta,\delta>0$ so that the two strips and $[0,\delta N)$ together account for less than $\varepsilon N/2$ eventually. Any large $N$ with $D_N(x)>\varepsilon$ must then have a bulk nonrectangular plateau $[U,V)$ with $\den(R_n)\le M$ meeting $[\delta N,N)$. Since $V/U\le p$, each such plateau covers at most $O_{p,\delta}(1)$ integer logarithmic bins $[m,m+1)$. For each fixed bulk shape, the logarithmic starts $s_j=\log U_j$ satisfy $s_{j+1}\ge c\exp(s_j)$ for a constant $c=c(B,p,q,M)>0$; once $c\exp(s_j)>s_j+\log(1/c)$, two steps satisfy $s_{j+2}\ge c\exp(c\exp(s_j))>\exp(s_j)$. Counting bins up to $T$ over the finitely many bulk shapes yields $O_{B,p,q,\varepsilon}(\log^* T)$. Summing $O(1/m)$ and $O(1)$ over the bad bins yields both displayed estimates in (ii).

(iii) From (ii), $\sum_{N\le X}D_N(x)/N=o(\log X)$, which forces $D^-(x)=0$. For any arc $J$, put $E_N=\sum_{n=0}^{N-1}(\mathbf1_J(\{B^n x\})-|J|)$, so $|E_N|\le ND_N(x)$. Summation by parts gives
\[
 \left|\sum_{n=0}^{X-1}\frac{\mathbf1_J(\{B^n x\})-|J|}{n+1}\right|
 =\left|\frac{E_X}{X}+\sum_{N=1}^{X-1}\frac{E_N}{N(N+1)}\right|
 \le 1+\sum_{N=1}^{X}\frac{D_N(x)}{N}=o(\log X),
\]
uniformly in $J$. Dividing by the harmonic sum $H_X\sim\log X$ proves logarithmic equidistribution (in the sense of Tsuji \cite{tsuji}).
\end{proof}

\section{The geometry of a single shape and certified bounds}
For fixed positive $a,d$ at sufficiently late shells ($L>a,K>d$), put $r=Q/P$. The plateau of shape $(a,d)$ consists of the integers in $[PU,PV)$ when $U<V$, where
\[
 U=\max(1+r/q^a,r+1/p^d),\qquad
 V=\min(1+r/q^{a-1},r+1/p^{d-1},p,qr).
\]
These inequalities express the last included and first omitted insertions on both top axes, clipped to the shell $[P,pP)\cap[Q,qQ)$. (When $a=L$ or $d=K$, the corresponding empty-fibre lower bound shifts by at most one integer position.) The continuous relative length of $[PU,PV)$ through any cutoff is maximized at $PV$ with value $1-U/V$, so the integer plateau weight is at most $1-U/V+O(1/P)$.

\begin{proposition}\label{th:shape}\evidence{Lean-supported paper proof}
The continuous relative length $1-U/V$ of the primary $(1,1)$ plateau has exact supremum
\[
 M_{p,q}=\frac{(p-1)(q-1)}{2pq-p-q},
\]
uniquely attained in $r$ at $r_*=q(p-1)/(p(q-1))$. For $a\ge2$ the relative length is strictly below $(q-1)/q^2$, and for $d\ge2$ strictly below $(p-1)/p^2$; both are strictly smaller than $M_{p,q}$. Under the simple-zero conditions, every fixed joint-depth failure set has ordinary upper density at least $M_{p,q}$.
\end{proposition}
\begin{proof}
For $(a,d)=(1,1)$ we have $V\le1+r$, giving $1-U/V\le\min((1-1/q)r,1-1/p)/(1+r)$. The two branches increase and decrease respectively, meeting at $r_*$ with common value $M_{p,q}$. Since $(p-1)(q-1)>1$, we have $1/(q-1)<r_*<p-1$, so $1+r_*<\min(p,qr_*)$ and shell clipping is inactive at $r_*$. The \lean{Phase17DiscrepancyBounds}{184}{general real inequality} and its \lean{Phase17DiscrepancyBounds}{232}{equality at the crossing} are kernel-verified (see also \lean{Phase16StonehamBounds}{232}{the binary case $4/11$}).

For $a\ge2$, feasibility $U<V$ gives $r<1+r/q^{a-1}$, so $y=r/q^{a-1}<1/(q-1)$. Comparing the first entries of $U$ and $V$ gives $1-U/V\le(1-1/q)y/(1+y)<(q-1)/q^2$. Exchanging the axes gives $(p-1)/p^2$ for $d\ge2$. Since $M_{p,q}^{-1}=p/(p-1)+q/(q-1)\le3/2+5/4=11/4$, we have $M_{p,q}\ge4/11>2/9\ge\max((p-1)/p^2,(q-1)/q^2)$.

Finally, for fixed congruence classes $K\equiv K_0\bmod A_0$ and $L\equiv L_0\bmod D_0$, the fractional parts of $iA_0\log p/(D_0\log q)$ are dense in $[0,1)$ by irrationality of $\log p/\log q$, so $r=q^L/p^K$ approaches $r_*$ along arbitrarily late shells in those classes. Under the simple-zero conditions, imposing any prescribed pair of fixed loss depths on shape $(1,1)$ and taking cutoffs at the right plateau endpoints proves that the ordinary upper density of every fixed joint-depth failure set is at least $M_{p,q}$.
\end{proof}

\begin{proposition}\label{th:cert}\evidence{Finite certificate}
For each row in the table below, the ordered axes $(p,q)$ and $(q,p)$ satisfy the simple-zero conditions, and in the shape grid $[1,A]\times[1,D]$ no two distinct shapes have both of their lifted loss residues within the radii $(R_L,R_K)$ in circular distance modulo $(up^{h-1},vq^{k-1})$:
\[
\begin{array}{c|c|c|c|c}
(B,p,q)&(\eta,\delta)&(A,D)&(R_K,R_L)&(h,k)\\ \hline
(2,3,5)&(10^{-10},10^{-11})&(15,21)&(25,17)&(9,7)\\
(2,3,7)&(10^{-6},10^{-7})&(8,13)&(16,10)&(3,9)\\
(2,3,13)&(10^{-6},10^{-7})&(6,13)&(16,8)&(2,7)\\
(36,5,7)&(10^{-6},10^{-7})&(8,9)&(12,10)&(5,4)
\end{array}
\]
Separately, for $(B,p,q)=(2,3,5)$, every nonempty shifted prefix $R_n$ through $n=3^{12}=531441$ has reduced denominator at least $15$ (attained), so no nonempty integer prefix occurs in this range.
\end{proposition}
\begin{proof}
The periods $(u,v)$ are respectively $(2,4)$, $(3,6)$, $(3,6)$, and $(4,6)$; the initial block sums modulo $(p^2,q^2)$ are $(6,15)$, $(6,7)$, $(6,65)$, and $(10,21)$, with first-layer zeros only at the block ends. For each shape, the root $\zeta$ is lifted one $p$-adic digit at a time by evaluating $G/p$ via the truncated Mahler--Newton sum of Proposition~\ref{th:local} (which is exact modulo $p^h$ at index $k=h-2$) and converting to $a+uz\bmod up^{h-1}$ and $d+vz'\bmod vq^{k-1}$. For every unordered pair of distinct shapes, comparing the circular distances against $(R_L,R_K)$ yields $0$ collisions in all four rows. The moduli $(up^{h-1},vq^{k-1})$ are $(13122,62500)$, $(27,34588806)$, $(9,28960854)$, and $(2500,2058)$, respectively. For the prefix scan, evaluating $\den(R_n)$ at every insertion or shell boundary and its predecessor through $n=3^{12}$ (between which $\den(R_n)$ is constant) gives minimum $15$. Exact integer scripts reproduce the \repo{scripts/phase16_mixed_product.py}{first grid} and the \repo{scripts/phase17_mixed_product_orbits.py}{remaining grids and prefix scan}.
\end{proof}

\begin{corollary}\label{th:examples}\evidence{Lean-supported paper proof}
For each row of Proposition~\ref{th:cert} and every fixed $M\ge1$, both the upper density of $\{n:\den(R_n)\le M\}$ and $D^+(\alpha(B,p)\alpha(B,q))$ are at most
\[
 S=M_{p,q}+\left(\frac p{p-1}+\frac q{q-1}\right)\eta+\delta .
\]
In row order, the bounds are $4/11+2.85\cdot10^{-10}$, $11250083/30000000$, $720004991/1860000000$, and $720004379/1740000000$.
\end{corollary}
\begin{proof}
Discard $[0,\delta N)$ and the two boundary strips of Theorem~\ref{th:sparse}. In the remaining window $[\delta N,N)$, any two positions have shell-index differences bounded in absolute value by $R_K,R_L$, since $p^{R_K-1},q^{R_L-1}>1/\delta$, and every bulk nonrectangular shape lies in $[1,A]\times[1,D]$ since $q^A,p^D>1/\eta$. For sufficiently large $N$, every position in this window with $\den(R_n)\le M$ is nonrectangular (by rectangular denominator growth in Proposition~\ref{th:local}) and loses at least $(h,k)$ layers. Proposition~\ref{th:cert} permits at most one bulk shape, and since $up^{h-1}>R_L$ and $vq^{k-1}>R_K$, even that single shape occurs in at most one shell. Proposition~\ref{th:shape} bounds its weight by $M_{p,q}+O(1/N)$, yielding upper density at most $S$ for $\{n:\den(R_n)\le M\}$. Hence every $w\in\mathcal W$ satisfies $\sum_{m_c\le M} w_c\le S$ for every $M$, so $\sum_c w_c\le S$. Theorem~\ref{th:mixture} and convexity of $\Delta$ give $D^+(x)\le\max_{w\in\mathcal W}\sum_c w_c\le S$.
\end{proof}

\section{Integer prefixes at the threshold}\label{sec:threshold}
\begin{proof}[Completion of the proof of Theorem~\ref{th:binary}]
Write $m=3^r5^s>1$ for a reduced denominator of a nontrivial binary cycle and let $m_0\in\{3,5,15\}$ be its radical. Since $v_3(2^2-1)=1$ and $v_5(2^4-1)=1$, lifting the exponent gives $\operatorname{ord}_{3^r}(2)=2\cdot3^{r-1}=\varphi(3^r)$ ($r\ge1$) and $\operatorname{ord}_{5^s}(2)=4\cdot5^{s-1}=\varphi(5^s)$ ($s\ge1$). Taking the least common multiple for $r,s\ge1$ gives $\operatorname{ord}_m(2)=4\cdot3^{r-1}5^{s-1}=\frac12\varphi(m)$. In every case the index of $\langle2\rangle$ in $(\Z/m\Z)^\times$ equals its index in $(\Z/m_0\Z)^\times$, so each cycle consists of all lifts modulo $m$ of its numerator coset modulo $m_0$: the full unit groups at $m_0=3,5$, and either $\{1,2,4,8\}$ or its negative $\{7,11,13,14\}$ at $m_0=15$.

For any probability $\mu$ on $\T$, interval errors are differences of the centered distribution function $F_\mu(t)=\mu([0,t))-t$, so $\Delta(\mu)=\sup F_\mu-\inf F_\mu$ (including one-sided limits). For a finite atomic cycle, $F_\mu$ decreases linearly between atoms, so its extrema occur immediately before and after atoms. Direct evaluation at $m_0=3,5,15$ gives $\Delta(\nu_c)=2/3,2/5,11/20$, respectively (the two $15$-cycles are reflections). Repetition of the same pattern in $m/m_0$ equal subintervals scales the range by $m_0/m$. (For illustration, distinct cycles need not add: assigning weight $1/4$ to each $15$-cycle and $1/2$ to $\lambda$ gives discrepancy $7/60$ rather than $11/40$.) In particular, every nontrivial cycle $c\ne0$ satisfies $\Delta(\nu_c)\le2/3$ and has period at least $4m/15$.

Put $s_*=4/11$ and $\gamma=2.85\cdot10^{-10}$. For any $w\in\mathcal W$, let $w_0$ be the weight of the zero cycle $\delta_0$ and $w'=\sum_{c\ne0} w_c$. By Corollary~\ref{th:examples}, $w_0+w'\le s_*+\gamma$. Convexity gives
\[
 \Delta(\sigma(w))\le w_0+\frac23 w'=(w_0+w')-\frac13 w'\le s_*+\gamma-\frac13 w'.
\]
Thus if $\Delta(\sigma(w))>s_*$, then $w'<3\gamma$ and $w_0>s_*-2\gamma$, which implies $w_0\ge\frac{11}4 w'$.

For any $t\in(0,1)$ and any nontrivial cycle $c$ of denominator $m>1$, the cycle points are distinct nonzero grid fractions $j/m$ ($1\le j\le m-1$). There are $\lceil mt\rceil-1<mt$ such points in $[0,t)$ and $m-\lceil mt\rceil\le m(1-t)$ in $[t,1)$. Since each atom has mass at most $15/(4m)$, we obtain $\nu_c([0,t))\le\frac{15}{4}t$ and $\nu_c([t,1))\le\frac{15}{4}(1-t)$, hence
\[
 -\frac{11}4(1-t)\le \nu_c([0,t))-t\le \frac{11}4 t.
\]
Summing against $w_c$ and adding $w_0(1-t)$ gives, for $0<t<1$,
\[
 \Bigl(w_0-\frac{11}4 w'\Bigr)(1-t)\le F_{\sigma(w)}(t)\le w_0+\Bigl(\frac{11}4 w'-w_0\Bigr)t.
\]
Under $w_0\ge\frac{11}4 w'$, every value and one-sided limit of $F_{\sigma(w)}$ lies in $[0,w_0]$, while $F_{\sigma(w)}(0)=0$ and $F_{\sigma(w)}(0^+)=w_0$, so $\Delta(\sigma(w))=w_0$.

For every $w\in\mathcal W$, shrinking arcs $[0,\delta)$ around $0$ gives $\Delta(\sigma(w))\ge w_0$. Because $w_{0,N}=\frac1N\#\{0\le n<N:R_n\in\Z\}$, diagonal extraction yields $\max_{w\in\mathcal W} w_0=C_0$. By Theorem~\ref{th:mixture}, $D^+(x)\ge C_0$; and whenever $D^+(x)>s_*$, a maximizing profile $w^*$ has $D^+(x)=\Delta(\sigma(w^*))=w_0^*\le C_0$, forcing $D^+(x)=C_0$. Hence $D^+(x)\le4/11\iff C_0\le4/11$, and if either fails then $D^+(x)=C_0$. The upper bound $D^+(x)\le s_*+\gamma$, $D^-(x)=0$, and logarithmic equidistribution follow from Corollary~\ref{th:examples} and Theorem~\ref{th:sparse}.
\end{proof}

\begin{remark}
The finite prefix scan in Proposition~\ref{th:cert} shows that no nonempty integer prefix occurs for $\alpha(2,3)\alpha(2,5)$ through $n=3^{12}=531441$. That finite computation is independent of the asymptotic equivalence in Theorem~\ref{th:binary}; whether any nonempty integer prefix exists beyond $3^{12}$, whether $C_0=0$, and whether $\alpha(2,3)\alpha(2,5)$ is binary normal all remain open.
\end{remark}

\section{Isolated insertions for independent integer parameters}
Before applying higher-precision $p$-adic lifting, elementary isolated-insertion windows already establish irrationality, transcendence, and discrepancy strictly below $1$ for arbitrary multiplicatively independent integer parameters. These results were previously placed in the first version of the companion Stoneham-polynomial manuscript and are consolidated here.

\begin{theorem}\label{th:43}\evidence{Kernel-verified}
The binary product $\alpha(2,3)\alpha(2,5)$ is irrational and satisfies $D^+(\alpha(2,3)\alpha(2,5))\le 43/45$.
\end{theorem}
\begin{proof}
See \lean{ProductThreeFive}{77}{\texttt{ProductThreeFive.discrepancy}} and \lean{ProductThreeFive}{93}{\texttt{ProductThreeFive.irrational}}. For each large $P=3^k$, choose $Q=5^\ell\in[P/3,5P/3)$ and set $s=P+Q$, $H=\lfloor4P/15\rfloor$. All competing sums at most $P,Q$ are at most $\max(P+Q/5,P/3+Q)$, and all sums with a larger component are at least $\min(3P,5Q)$, so both gaps around $s$ are at least $4P/15$ and $[s-H,s+H)\subset[P,3P)$ ends by $44P/15$. Since $R_s-2^H R_{s-H}=1/(PQ)$, at least one of the two adjoining windows of length $H$ retains denominator divisible by $P$. The fixed-support orbit estimate and vanishing shifted tail give uniform arc error $o(P)$ there. For $238R\le N<714R$ ($R=3^h$), the four complete disjoint windows from $P=3R,9R,27R,81R$ have combined length at least $32R-4\ge2N/45$ eventually, yielding $D^+\le43/45$. Irrationality follows because uniformity on these windows is impossible for a rational orbit confined to a finite grid.
\end{proof}

\begin{remark}
Between Theorem~\ref{th:43} and the multi-depth bound of Theorem~\ref{th:binary}, three progressively tighter upper bounds on $D^+(\alpha(2,3)\alpha(2,5))$ are also kernel-verified at the pinned commit: $27/40$ by counting all surviving top quinary columns (\lean{ProductThreeFiveCoverage}{176}{\texttt{column\_discrepancy}}), $143881/250000=0.575524$ by first-precision phase coverage (\lean{ProductThreeFiveOrdinary}{268}{\texttt{mixed\_product\_discrepancy}}), and $4631/12500=0.37048$ by removing full rectangles via \lean{ProductThreeFiveRectangleValuation}{336}{rectangular denominator growth} (\lean{ProductThreeFiveNonrectangularOrdinary}{246}{\texttt{mixed\_product\_discrepancy}}).
\end{remark}

\begin{theorem}\label{th:indep}\evidence{Lean-supported paper proof}
If $c,d\ge2$ are multiplicatively independent integers with $\gcd(B,cd)=1$, then $y=\alpha(B,c)\alpha(B,d)$ is transcendental and $D^+(y)<1$.
\end{theorem}
\begin{proof}
Since $(c-1)(d-1)>1$, the open interval $(1/(d-1),c-1)$ is nonempty. Choose a compact subinterval $[r_1,r_2]$ of positive length whose base-$d$ logarithmic length is smaller than $1$. Because $\log_d c$ is irrational, the backward translates of the corresponding open arc cover $\T$, so compactness gives a set of exponents $k$ with bounded gaps $H_0$ such that for each sufficiently large $k$ in this set ($c^k r_1\ge d$), some $\ell\ge1$ has $r=d^\ell/c^k\in[r_1,r_2]$. Put $P=c^k$, $Q=d^\ell$, and $s=P+Q$. The lower gap from competing rectangle sums is at least $\min((1-1/d)Q,(1-1/c)P)$, and the upper gap from sums outside the rectangle is at least $\min((c-1)P-Q,(d-1)Q-P)$; both are uniformly $\gg P$. Choose $\eta>0$ small enough that $[s-\lfloor\eta P\rfloor,s+\lfloor\eta P\rfloor)\subset[P,cP)$ contains no insertion other than $s$.

Fix a prime $\pi\mid c$. The jump $1/(PQ)$ has $\pi$-adic pole order at least $kv_\pi(c)$, and $B$ is a $\pi$-adic unit, so at least one of the two adjacent windows of length $\lfloor\eta P\rfloor$ retains reduced denominator at least $\pi^{kv_\pi(c)}$. Its denominator is supported on the prime divisors of $cd$ and is at most $n^2$. The \lean{SmoothOrbit}{207}{fixed-support orbit estimate} and vanishing tail give uniform arc error $o(P)$ on that window (also when $c$ and $d$ share primes). For $c^K\le N<c^{K+1}$, an eligible shell with exponent in $[K-H_0,K-1]$ is complete and supplies a uniform window of length at least $\eta N/(2c^{H_0+1})$ eventually, giving $D^+(y)\le1-\eta/(2c^{H_0+1})<1$ and irrationality.

For transcendence, truncate the unshifted series at $s$. Its rational value $u/t$ (in lowest terms) has denominator $t$ supported on the prime divisors of $Bcd$ with $t\le B^s s^2$. Every omitted term has exponent at least $s+\eta P$, and the coefficient weights are summable, so $0<y-u/t=O(B^{-s-\eta P})$. Since $s/P=1+r\le1+r_2$ is bounded, there is $\epsilon>0$ such that $0<y-u/t<t^{-1-\epsilon}$ for all sufficiently late selected insertions. Positivity and irrationality of $y$ produce infinitely many distinct reduced approximants $u/t$. Ridout's restricted-denominator theorem \cite[Theorem, p.~125]{ridout57} (taking $\mu=1,\nu=0$; see also \cite{ridout}) excludes infinitely many such approximations to any real algebraic irrational, proving that $y$ is transcendental.
\end{proof}

\section{Quadratic plateaus on the dyadic axis}\label{sec:quad}
Now let $b\ge3$ be odd, $\alpha=\alpha(b,2)$, and $y=(A\alpha^2+E\alpha+C)/D$ for integers $A\ne0,E,C$ and $D\ge1$. Here let $R_n$ denote the shifted rational prefix of the \emph{numerator}, retaining ordered pairs $2^i+2^j\le n$ in the square, terms $2^j\le n$ in the linear sum, and $Cb^n$. Write $n_k=3\cdot2^{k-1}$ and $v_k=2^{k+1}$ for $k\ge1$; on $[n_k,v_k)$ the square contains the $k\times k$ rectangle except $(k,k)$, and $\den(R_n)=q_k$ is constant. Let $s_D\ge0$ be least with $\den(b^{s_D}/D)$ coprime to $b$, put $y_k=\{b^{s_D}R_{n_k}/D\}$ (which is purely periodic under multiplication by $b$), let $\Delta_b(y_k)$ be the extreme discrepancy of the uniform measure on its $b$-cycle, and put $w_{k,N}=|[0,N)\cap[n_k,v_k)|/N$. Set $s=v_2(b^2-1)-1\ge2$ and $[z]_+=\max(0,z)$.

\begin{theorem}\label{th:plat}\evidence{Kernel-verified}
For every such rational quadratic $y=(A\alpha^2+E\alpha+C)/D$,
\[
 D_N(y)-\max_{k\ge1} w_{k,N}\Delta_b(y_k)\longrightarrow0,\qquad
 D^+(y)=\frac14\limsup_{k\to\infty}\Delta_b(y_k),\qquad
 D^-(y)=0.
\]
Moreover, $\{b^n y\}$ is logarithmically equidistributed on $\T$, and $y$ is normal in base $b$ if and only if $q_k=\den(R_{n_k})\to\infty$. Writing $\delta_j=(b^{-2^j}-1)/2^j\in\Q_2$, $\lambda=\lim_{j\to\infty}\delta_j$, $C_{\mathrm{int}}=\lambda+\sum_{j\ge1}(\delta_j-\lambda)$, and $\kappa=-(2A+E+2AC_{\mathrm{int}})/(2A\lambda)\in\Q_2$, one has
\[
 q_k=2^{[k-v_2(2A)-s-v_2(k-\kappa)]_+}\qquad(k\ge1).
\]
In particular, if $\kappa$ is algebraic over $\Q$, then $y$ is normal in base $b$.
\end{theorem}
\begin{proof}
All assertions are kernel-verified with index $j=k-1$: the \lean{QuadraticPlateauSpectrum}{116}{all-prefix plateau theorem}, the \lean{QuadraticPlateauSpectrum}{167}{exact upper limit}, the \lean{ParameterTwoQuadraticOrbits}{33}{lower limit $D^-=0$}, \lean{ParameterTwoQuadraticOrbits}{42}{logarithmic equidistribution}, the \lean{ParameterTwoQuadraticOrbits}{16}{denominator normality criterion}, the \lean{ParameterTwoQuadratic}{187}{exact upper-shell denominator formula}, and \lean{ParameterTwoQuadraticOrbits}{22}{normality for algebraic $\kappa$}. Briefly, lower half-shells $[2^k,n_k)$ have growing denominator $2^{2k-2-v_2(A)}$ eventually and uniform counts, while upper half-shells $[n_k,v_k)$ contribute $w_{k,N}\nu_{y_k}$ up to $o(1)$. Because distinct shells with $q_i,q_j\le2^h$ satisfy $2^{i-v_2(2A)-s-h}\mid(j-i)$, \lean{ParameterTwoQuadraticHeight}{153}{at most one bounded denominator survives} in any sufficiently late fixed window of shell indices; all large denominators give small cycle error by the fixed-support estimate, isolating the single maximum whose weight peaks at $N=v_k$ with value $1/4$.
\end{proof}

When the radix $b$ varies at a fixed shell $k$ for $X^2-2X$, the resonant radices are also completely classified by a kernel-verified $2$-adic isometry.

\begin{proposition}\label{th:resrad}\evidence{Kernel-verified}
Take $A=1,E=-2$ in Theorem~\ref{th:plat}. Fix integers $k\ge1$, $s_0\ge2$, and $H\ge0$ such that $t=k-1-s_0-H>s_0$. Among odd integer radices $b\ge3$ with $v_2(b^2-1)=s_0+1$:
\begin{itemize}[leftmargin=2em,itemsep=2pt]
\item if $v_2(k+1)\ne s_0$, no such radix satisfies $q_k(b)\le2^H$;
\item if $v_2(k+1)=s_0$, the radices satisfying $q_k(b)\le2^H$ form exactly two opposite residue classes modulo $2^t$.
\end{itemize}
\end{proposition}
\begin{proof}
This is formalized in \lean{ResonantRadixClasses}{69}{\texttt{target\_valuation\_of\_hit}} and \lean{ResonantRadixClasses}{114}{\texttt{resonant\_residue}} (using the \lean{ResonantRadix}{199}{radix isometry} $b\mapsto\kappa(b)+1$ on $1+4\Z_2$). For example, at $k=11,s_0=2,H=0$ ($t=8$), the upper prefix is integral precisely for $b\equiv93,163\pmod{256}$.
\end{proof}

Returning to a fixed odd radix $b\ge3$, evaluating the dyadic cycle discrepancies $\Delta_b(y_k)$ explicitly sharpens the general $1/4$ ceiling of Theorem~\ref{th:plat}.

\begin{proposition}\label{th:dyadic}\evidence{Lean-supported paper proof}
Suppose $D=1$ and write $y_k=u_k/2^{h_k}$ in lowest terms modulo $1$, with $h_k=0$ when $y_k=0$. Put $H=\liminf_{k\to\infty} h_k\in\N\cup\{\infty\}$.
\begin{enumerate}[label=\textup{(\roman*)},leftmargin=2em,itemsep=2pt]
\item If $b\equiv1\pmod4$, then $D^+(y)=2^{-2-[H-s]_+}$ (with $2^{-\infty}=0$).
\item For every odd $b$, $D^+(y)=1/4$ if and only if $h_k\le v_2(b-1)$ infinitely often.
\item If $b\equiv3\pmod4$ and $H\ge2$, then $D^+(y)\le\frac14(1-2^{-s})\,2^{-[H-s-1]_+}$.
\item For $P(X)=X^2-2X+C$, define integers $Z_1=2-2b^2$ and $Z_{k+1}=2+2b^{2^k}(Z_k-1)$. Then for all $k\ge1$,
\[
 y_k=\left\{-\frac{Z_k b^{2^{k-1}}}{2^{2k-1}}\right\},\qquad
 h_k=[2k-1-v_2(Z_k)]_+ .
\]
\end{enumerate}
\end{proposition}
\begin{proof}
For odd $u$ (or $u=0,h=0$), the $b$-orbit of $u/2^h$ is a singleton precisely when $h\le v_2(b-1)$. If $b\equiv1\pmod4$, then $v_2(b-1)=s$, and repeated squaring gives $v_2(b^{2^j}-1)=s+j$, so $\langle b\rangle\bmod 2^h$ is $\{1+2^s m\bmod 2^h\}$. The orbit of $u/2^h$ is a translated uniform grid of spacing $g=2^{-[h-s]_+}$, whose extreme discrepancy is $g$. Taking $\limsup_k$ in Theorem~\ref{th:plat} proves (i).

If $b\equiv3\pmod4$, then $v_2(b-1)=1$ and $v_2(b^2-1)=s+1$, so $\langle b^2\rangle\bmod 2^h$ is $\{1+2^{s+1}m\bmod 2^h\}$. For $h\ge2$, the two cosets of $\langle b^2\rangle$ give two equal grids of spacing $g=2^{-[h-s-1]_+}$ separated by offset $\theta g$, where $\theta=\{(b-1)u/2^{\min(h,s+1)}\}$. In each cell of width $g$, the centered cumulative function has two equal jumps of size $g/2$ followed by descents of lengths $\theta g$ and $(1-\theta)g$, giving range $g\max(\theta,1-\theta)$. Since $v_2(b-1)=1$, the reduced denominator of $\theta$ divides $2^s$, so $\min(\theta,1-\theta)\ge2^{-s}$ and $\Delta_b(u/2^h)\le(1-2^{-s})g$. Thus every nonsingleton dyadic cycle has discrepancy uniformly bounded away from $1$ (by $1/2$ if $b\equiv1\bmod4$ and by $1-2^{-s}$ if $b\equiv3\bmod4$). Passing to $\limsup_k$ in Theorem~\ref{th:plat} gives (ii) and (iii).

For (iv), write $l_j=2^{-j}b^{-2^j}=2^{-j}+\delta_j$, $S_k=\sum_{j=1}^k l_j$, and $J=1+\sum_{j=1}^k\delta_j$, so $S_k=J-2^{-k}$ with $J,\delta_k\in\Z_2$. The rectangle-minus-corner identity gives
\[
 b^{-n_k}R_{n_k}=S_k^2-l_k^2-2S_k+C
 =J^2-\delta_k^2-2J+C-2^{1-k}\Bigl(\sum_{j=1}^k\delta_j+\delta_k\Bigr).
\]
By the \lean{Phase16StonehamBounds}{68}{cancellation identity} (with index $j=k-1$), $\sum_{j=1}^k\delta_j+\delta_k=Z_k/(2^k b^{2^k})$. Hence $R_{n_k}+Z_k b^{2^{k-1}}/2^{2k-1}$ belongs to $\Z[1/2]\cap\Z_2=\Z$, proving the formula for $y_k$. Reducing modulo $1$ gives $h_k$, matching the \lean{Phase16StonehamBounds}{126}{formal upper-shell denominator theorem}.
\end{proof}

\begin{corollary}\label{th:quarter}\evidence{Lean-supported paper proof}
For every odd radix $b\equiv2\pmod3$ and integer $C\not\equiv1\pmod3$,
\[
 D^+\!\left(\frac{\alpha^2-2\alpha+C}{3}\right)
 \le \frac14\left(1-\frac{1}{3\cdot2^{v_2(b+1)}}\right)<\frac14.
\]
For $b=5$, one has $D^+((\alpha^2-2\alpha)/9)\le5/72$.
\end{corollary}
\begin{proof}
Since $b\equiv2\pmod3$, $b^{2^j}\equiv1\pmod3$ and $l_j\equiv(-1)^j\pmod3$ for all $j\ge1$. Thus $S_k\equiv-1\pmod3$ for odd $k$ and $0\pmod3$ for even $k$. In both cases $S_k^2-l_k^2-2S_k+C\equiv C-1\not\equiv0\pmod3$, and multiplying by $b^{n_k}$ preserves nonvanishing modulo $3$. Therefore $y_k=\{R_{n_k}/3\}=\{c/3+u/2^h\}$ with $3\nmid c$ and $u$ odd (or $u=0$ if $h=0$). Set $t=v_2(b^2-1)\ge3$. For $0\le h\le t$, $b^2 y_k\equiv y_k\pmod1$ while $b y_k\not\equiv y_k\pmod1$, so the cycle is $\{y_k,b y_k\}$, whose discrepancy is $1-\|(b-1)y_k\|$ (where $\|\cdot\|$ is distance to $\Z$). Because $(b-1)y_k$ has nonzero third component $c/3$ and $2$-part of denominator dividing $2^{h-v_2(b-1)}\mid 2^{v_2(b+1)}$, we have $\|(b-1)y_k\|\ge1/(3\cdot2^{v_2(b+1)})$, giving $\Delta_b(y_k)\le1-1/(3\cdot2^{v_2(b+1)})$. For $h>t$, $\langle b^2\rangle$ fixes the third component and acts as $1+2^t\Z_2$ on the $2$-component, so the cycle is a union of two grids of spacing $2^{t-h}\le1/2$, whence $\Delta_b(y_k)\le1/2<1-1/(3\cdot2^{v_2(b+1)})$. Applying Theorem~\ref{th:plat} proves the first bound.

For $b=5,D=9,C=0$, the same reduction modulo $3$ gives $3\nmid R_{n_k}$, forcing a primitive ninth component $c/9$ ($c\in\{1,2,4,5,7,8\}$). Since $\operatorname{ord}_9(5)=6$ and $v_2(5^6-1)=3$, for $h\ge5$ the subgroup $\langle5^6\rangle$ yields a grid bound $2^{3-h}\le1/4$. For $0\le h\le4$, exact evaluation of the cumulative range over all admissible $(c,u)$ gives maxima $2/9,2/9,2/9,5/18,5/36$, respectively (reproduced by \repo{scripts/phase17_dyadic_stoneham.py}{the dyadic Stoneham checker}). Thus $\Delta_5(y_k)\le5/18$ for all $k$, and Theorem~\ref{th:plat} gives $D^+\le5/72$.
\end{proof}

For a general divisor $D=2^e D_1$ coprime to $b$ ($D_1$ odd), the $D_1$-component of $y_k$ is eventually periodic because $(S_k,2^{-k},b^{2^{k-1}})\bmod D_1$ evolves in a finite state space; with $l=\operatorname{lcm}(2,\operatorname{ord}_{D_1}(b))$, $\langle b^l\rangle$ gives the grid bound $2^{v_2(b^l-1)-h}$ for $h>v_2(b^l-1)$, reducing any cycle-supremum bound to finitely many small dyadic levels. Whether $q_k\to\infty$ (equivalently, whether $\alpha^2-2\alpha$ is base-$b$ normal) for any fixed odd integer radix $b\ge3$ remains open: in $2$-adic terms $\kappa_0(\lambda)=-C_{\mathrm{int}}(\lambda)/\lambda$ is transcendental as a function of $\lambda$ over $\C(\lambda,e^{\lambda/2^m})$, but standard $p$-adic linear-form and Mahler-type methods \cite{bugeaud_laurent,nishioka} do not decide $v_2(k-\kappa_0(b))\ge k-O(1)$ at a fixed integer radix $b$.

\section*{Formal verification and data availability}
All repository links are pinned to research-checkpoint-18 (\repo{}{commit \texttt{e4809cf}}). The \repo{notes/phase18/verification.md}{checkpoint verification record} documents the full Lean 4.36.0 build, axiom audit (accepting only \texttt{propext}, \texttt{Classical.choice}, and \texttt{Quot.sound}), and independent kernel replay. Complete kernel-verified statements in this manuscript are Theorem~\ref{th:43}, Theorem~\ref{th:plat}, and Proposition~\ref{th:resrad}. Theorems~\ref{th:binary}, \ref{th:mixture}, \ref{th:sparse}, and~\ref{th:indep}, Propositions~\ref{th:local}, \ref{th:shape}, and~\ref{th:dyadic}, and Corollaries~\ref{th:examples} and~\ref{th:quarter} are Lean-supported paper proofs whose formalized steps are linked inline; Proposition~\ref{th:eisen} is a paper proof; and Proposition~\ref{th:cert} is an exact finite certificate. Exact integer and rational scripts reproducing all finite certificates and cycle tables are available at \repo{scripts/phase16_mixed_product.py}{\texttt{phase16\_mixed\_product.py}}, \repo{scripts/phase17_mixed_product_orbits.py}{\texttt{phase17\_mixed\_product\_orbits.py}}, and \repo{scripts/phase17_dyadic_stoneham.py}{\texttt{phase17\_dyadic\_stoneham.py}}. The accompanying verification map (\texttt{VERIFICATION.md}) lists exact Lean declarations, external analytic inputs, and reproduction commands.

\section*{Use of AI tools}
This manuscript was produced in an AI-assisted research project. Large language models proposed conjectures and proof strategies, wrote the proofs and the text, implemented the exact computations and Lean 4 formalizations, searched the literature and reviewed drafts across several model providers. The repository owner directed the project and is responsible for its release. The Lean kernel checked the declarations listed in the verification map. The written proofs, literature comparisons and priority statements have been checked only by AI models from different providers. They have not been checked by an independent human expert. Cross-model review is not peer review.

\bibliographystyle{plain}
\bibliography{references}
\end{document}
